\section{Implementation Roadmap for Optional GPS Observations}
\label{sec:gps-observation-implementation-roadmap}

The mathematical contract for the optional GPS observation channel is given in
Section~\ref{sec:gps-trip-attribute-observations}.  The following phases record
the proposed implementation order.  A phase is marked implemented only after
its code, tests, and documentation have been reviewed; the remaining phases
are explicitly marked \emph{Not implemented}.  The absence of an auxiliary
data file must remain a supported
configuration throughout: an empty observation bundle must reproduce the
current count-only objective, gradients, fingerprints, checkpoints, and
reports.

\subsection{Phase 0: Aggregate data contract}
\label{sec:gps-implementation-phase-0}
\textbf{Status: Not implemented.}

Define the privacy-safe data product that can be prepared locally by the data
provider.  The minimum product is a versioned histogram of in-vehicle travel
time and transfer-count bins, optionally stratified into overall, morning-peak,
off-peak, and evening-peak periods.  Its metadata must contain units, bin
boundaries, stratum definitions, collection period, valid and excluded journey
counts, cleaning reasons, APC-overlap policy, and a file checksum.  It must also
provide either a Dirichlet concentration, an effective sample size, or an
explicit statement that only a tempered multinomial is justified.  Exact GPS
traces and personal identifiers are outside the package contract.

The phase is complete only when a small aggregate fixture and its metadata have
been approved.  The fixture must include at least one sparse stratum, one
excluded-data record, and the no-service-hour convention.

\subsection{Phase 1: Generic optional observation interface}
\label{sec:gps-implementation-phase-1}
\textbf{Status: Implemented (generic interface; 2026-08-20).}

Introduce a generic gravity observation-channel abstraction.  A channel exposes
its observed aggregate data, a prediction from gravity demand, a JAX-compatible
likelihood contribution, a stable data fingerprint, and a serializable report
breakdown.  The gravity objective accepts a tuple of channels, defaulting to an
empty tuple.  The existing APC count likelihood remains the primary built-in
channel and is not rewritten as a GPS-specific special case.

The public implementation is provided by
\path{public_transportation.inference.gravity.observations}.  It validates
channel names, kinds, fingerprints, report mappings, and the free-OD dimension
at objective construction.  An empty bundle is normalized as the default and
does not alter the count-only objective or its model fingerprint; an enabled
bundle contributes its channel identities to the model fingerprint.  No
provider-specific data loader, trip-attribute response operator, or auxiliary
likelihood is enabled by this phase.

The required regression test is exact count-only equivalence: constructing the
objective without channels must preserve the current scalar objective,
gradient, model fingerprint, checkpoint identity, and result fields.

\subsection{Phase 2: Aggregate loader and validation}
\label{sec:gps-implementation-phase-2}
\textbf{Status: Implemented (loader and validation; 2026-08-20).}

Add a loader for the provider's aggregate histogram and metadata.  Validation
must reject duplicate bins, negative or non-integral counts, inconsistent
stratum totals, overlapping or incomplete bin definitions, invalid units,
missing checksums, and unsupported likelihood settings.  An absent optional
file must be interpreted as ``no auxiliary observations''; an explicitly
enabled but missing or invalid file must fail before expensive operator
construction.

The validator must write a compact audit containing the source fingerprint,
valid and excluded counts, stratum totals, bin definitions, overlap policy, and
uncertainty metadata.  Tests must cover both the absent-data path and malformed
aggregate inputs.

The public implementation is
\path{public_transportation.inference.gravity.aggregate}.  It accepts a
versioned JSON document containing one or more numeric histograms.  Each
histogram declares its attribute, unit, finite support, contiguous closed-open
bin boundaries, and named strata.  Stratum counts may be supplied as an array
in bin order or as an object keyed by bin label; in both cases the declared
total must equal the sum of non-negative integer counts.  The document also
requires collection and cleaning metadata, an APC-overlap policy, an
uncertainty declaration (multinomial, Dirichlet--multinomial with a
concentration/effective sample size, or tempered multinomial), and a semantic
SHA-256 content checksum.  The loader records the actual file checksum and a
path-independent aggregate fingerprint in its audit.

When the optional channel is disabled, the loader returns no observation and
records an ``absent`` audit without reading the file.  When it is enabled, a
missing or invalid file raises before routing or operator construction and
records a ``rejected`` audit.  A valid file records a ``validated`` audit with
the source path, both checksums, histogram and stratum summaries, metadata, and
uncertainty settings.  This phase does not yet build the route-response
operator or add the aggregate likelihood; those are deferred to Phases 3--5.

\subsection{Phase 3: Trip-attribute response operator}
\label{sec:gps-implementation-phase-3}
\textbf{Status: Implemented (reference and sharded response operator; 2026-08-20).}

Extend routing preparation only when the auxiliary channel is enabled.  For
each free OD--time cell, persist route-share-weighted contributions to the
declared attribute bins and departure strata.  The result is the operator
$B$ in equation~\eqref{eq:gps-attribute-mass}, together with its fixed-demand
offset.  It must use the same route choices, compact OD layout, timetable
identity, and feasibility contract as the APC operator $A$.

The implementation must not infer a distribution by binning only the mean
journey-time feature when several assigned paths exist.  It must preserve the
path-share distribution, using a compact or sharded matrix-free representation
for large layouts.  Attribute-response artifacts receive a separate
fingerprint; changing GPS bins must not rebuild $A$.

Tests must compare forward and adjoint products on a small multi-route fixture,
include fixed positive demand, and verify that unsupported observed bins stop
preflight with a clear audit.

The public reference implementation is
\path{public_transportation.inference.gravity.attribute_operator}.  A
\texttt{GravityRouteShare} records the stratum, attribute-bin label, and
probability assigned to one route for one free OD--time cell.  The constructor
adds all route shares belonging to a category, so a cell with several routes
retains its categorical distribution rather than being represented by a mean
journey attribute.  The resulting
\texttt{GravityAttributeResponseOperator} exposes JAX-compatible forward and
adjoint products and a fixed-demand attribute offset.  It can store one dense
matrix for small fixtures or contiguous category shards for large artifacts;
the sharded products do not concatenate the complete response matrix merely
to apply it.  Every artifact records the compact OD, assignment, graph,
timetable, and feasibility fingerprints and has an identity separate from
the APC operator.  The optional \texttt{save}/\texttt{load} artifact path
verifies this identity on reload.

The function \texttt{validate\_aggregate\_support} is the strict Phase-3
support audit.  It compares positive aggregate counts with categories that
have nonzero route support, returns a machine-readable report, and raises
\texttt{GravityAttributeSupportError} for a positive unsupported bin.  Zero
counts in a category with no route support are not treated as observations.
This is a preflight contract: no auxiliary likelihood or expensive estimator
work is started after the audit fails.  Integration with the case-study
routing-preparation adapter and the aggregate likelihood remains in Phases 4
and 5.

\subsection{Phase 4: Distribution likelihoods and uncertainty}
\label{sec:gps-implementation-phase-4}
\textbf{Status: Implemented (aggregate likelihood functions; 2026-08-20).}

Implement the aggregate likelihoods as independent channel components:

\begin{enumerate}
  \item ordinary multinomial for a representative independent sample;
  \item Dirichlet--multinomial when clustering, reconstruction uncertainty, or
    day-to-day variation can be summarized by a concentration parameter or
    effective sample size;
  \item a clearly named tempered multinomial composite likelihood when no
    uncertainty estimate is available.
\end{enumerate}

The Dirichlet--multinomial and tempering parameters must not be silently
combined to downweight the same evidence twice.  Tests must cover the
multinomial limit at large concentration, overdispersion at smaller
concentration, the zero-weight limit, finite-difference gradients, and
normalization within each stratum.

The public implementation is
\path{public_transportation.inference.gravity.likelihoods}.  It normalizes
predicted masses independently within every declared stratum and evaluates
all histograms as separate contributions.  The Phase-2
\texttt{effective\_sample\_size} is interpreted as the Dirichlet concentration
for this first implementation; this is a declared approximation, not an
assumption that the raw GPS row count is an exposure.  A zero-mass stratum
with zero observed count contributes zero, while a positive observed count in
an unsupported stratum must have been stopped by the Phase-3 support audit.
The tempered form multiplies only the ordinary multinomial log likelihood.
It is rejected when a concentration or effective sample size is supplied, so
the same uncertainty is not downweighted twice.  The functions return JAX
arrays and support automatic differentiation; they do not yet modify the
count-only gravity objective (Phase 5).

\subsection{Phase 5: Objective, estimator, and provenance integration}
\label{sec:gps-implementation-phase-5}
\textbf{Status: Implemented (objective, gradients, estimator provenance; 2026-08-20).}

Add the channel contributions to the existing gravity objective and both
gradient strategies without changing the count-only path.  The objective
breakdown must report, separately, the APC log-likelihood, each auxiliary
log-likelihood, the tempering or concentration settings, the prior penalty,
and the total MAP objective.  The result and checkpoint metadata must include
auxiliary data fingerprints and channel specifications only when a channel is
enabled.

Changing or adding an auxiliary data source must invalidate the corresponding
fit checkpoint.  With no source enabled, old checkpoints and model identities
must remain compatible.  Tests must verify objective composition, adjoint and
batched-forward agreement, resume rejection after a source change, and exact
backward compatibility when the channel tuple is empty.

The concrete bridge is
\path{public_transportation.inference.gravity.aggregate_channel}.
It creates one Phase-1 observation channel per validated histogram, adds the
fixed-demand attribute offset to its prediction, and invokes the Phase-4
likelihood.  \texttt{GravityObjectiveEvaluation} now retains the APC
log-likelihood, total auxiliary log-likelihood, and per-channel contributions;
the total data likelihood and MAP objective are their explicit sums.  Both
the batched-forward and adjoint gradient strategies include the auxiliary
chain rule through the demand vector, while the empty bundle follows the
previous count-only path.

Estimator results, run manifests, and enabled-channel checkpoints retain the
aggregate/channel reports and fingerprints.  Disabled channels are omitted
from these records.  The model fingerprint includes the channel identity, so
changing aggregate content, uncertainty settings, route-response artifacts,
or bins rejects a resume with a fingerprint mismatch.  Existing count-only
checkpoints remain compatible because their identity and serialized shape are
unchanged.

\subsection{Phase 6: Case configuration, adapter, and reporting}
\label{sec:gps-implementation-phase-6}
\textbf{Status: Not implemented.}

Add an optional TOML contract, preferably in
\path{config/auxiliary_observations.toml}, that declares sources, paths,
strata, bins, likelihood family, concentration or effective sample size, and
tempering.  The generic adapter must expose a documented hook for constructing
channels while retaining the existing case-owned responsibility for data
preparation.

The check and preflight stages must report whether auxiliary data are absent,
validated, or rejected.  Fit reports must retain the aggregate input checksum,
cleaning audit, APC-overlap decision, effective sample size or concentration,
and every channel's likelihood contribution.  The template must remain usable
without the optional configuration file.

\subsection{Phase 7: Case-study pilot and sensitivity analysis}
\label{sec:gps-implementation-phase-7}
\textbf{Status: Not implemented.}

Run the first pilot on a small or reduced case before the full network.  Compare
the count-only fit with the GPS-augmented fit using one overall distribution
and the approved broad strata.  Vary the Dirichlet concentration or effective
sample size, the tempering factor when used, bin boundaries, and the treatment
of the 22~November overlap.  Report APC fit, GPS distribution fit, parameter
changes, runtime, and sensitivity.

The GPS term may be promoted to the full case only if it improves or stabilizes
the declared diagnostics without creating unsupported-bin failures or allowing
the auxiliary sample to dominate the exhaustive APC evidence.  If the data
provider cannot supply a defensible aggregate product, this phase remains
disabled and the existing count-only MAP workflow is the valid result.
